\documentclass[11pt]{letter}
\usepackage[utf8]{inputenc}
\usepackage[T1]{fontenc}
\usepackage[margin=1in]{geometry}
\usepackage{hyperref}
\usepackage{siunitx}
\urlstyle{same}

\signature{Jonathan Mace\\Corresponding author}
\address{Jonathan Mace\\Microsoft Research\\Redmond, WA\\jonathanmace@microsoft.com}
\date{\today}

\begin{document}

\begin{letter}{%
Professor A.\ V.\ Metrikine\\
Editor-in-Chief\\
\textit{Journal of Sound and Vibration}}

\opening{Dear Professor Metrikine,}
\small

Please consider our manuscript, ``Fill-Dependent Resonance of the Human Urinary Bladder: Competing Geometry and Stiffening in a Fluid-Filled Viscoelastic Shell,'' for publication in the \textit{Journal of Sound and Vibration}. The paper models the bladder wall as a fluid-filled viscoelastic shell whose radius, thickness, internal pressure, and stiffness vary with fill volume. This yields a mechanistic explanation for a U-shaped frequency-volume relation, with the fundamental flexural mode reaching a minimum near \SI{222}{\milli\litre} as geometric softening during filling competes with exponential detrusor strain-stiffening. The analysis then shows that clinically relevant excitation lies not at resonance itself but in a sub-resonant forced-response regime driven by pelvic transmissibility in the \SIrange{4}{8}{\hertz} whole-body vibration band.

We believe the manuscript fits \textit{JSV} particularly well because the core advance is in structural dynamics and modal interpretation. The paper treats the bladder as a shell-vibration problem with fluid loading, prestress, and fill-dependent material properties, and it uses that framework to distinguish resonance from forced response in a way that should be broadly relevant to vibration transmission in compliant biological structures. The coupling analysis also shows that mechanical excitation transmitted through the pelvis is approximately 6.4 x 10$^3$ times stronger than airborne acoustic excitation, reinforcing the vibrational rather than purely acoustical nature of the problem. As part of our broader Browntone research programme, the paper extends the same shell-theory framework to a clinically recognisable organ-level application.

The manuscript is original, has not been published previously, is not under consideration elsewhere, and has been approved by both authors. Brian R. Mace's submission email is \texttt{b.mace@auckland.ac.nz}. We would be pleased for the paper to be considered at \textit{JSV} as the primary target venue, given its emphasis on modal analysis, shell mechanics, and vibration transmission.

\closing{Yours sincerely,}

\end{letter}
\end{document}
